\documentclass[10pt,a4paper]{amsart}
\usepackage[margin=19mm]{geometry}
\usepackage{amsmath,amssymb,mathtools}
\usepackage[scr=rsfs,cal=euler]{mathalfa}
\usepackage{xfrac}
\usepackage[unicode,hidelinks,bookmarks=false]{hyperref}
\newcommand{\ie}{i.e.\ }
\newcommand{\cf}{cf.\ }
\newcommand{\resp}{respectively\ }
\newcommand{\Subsection}{\subsection}
\providecommand{\nobreakdash}{\nobreak\hbox{-}}
\setlength{\emergencystretch}{2em}
\multlinegap0pt
\newtheorem{lemma}{Lemma}
\usepackage{cleveref}
\crefname{lemma}{Lemma}{Lemmas}
\newcommand{\R}{\mathbf{R}}
\newcommand{\dd}{\mathop{}\!\dletter}
\newcommand{\dletter}{\mathrm{d}}
\newcommand{\step}[1]{\smallskip\noindent\emph{#1}}
\title{Late-time tails for linear waves on radially symmetric stationary spacetimes of two space dimensions}
\author{Onyx Gautam}
\date{Source: arXiv:2605.03220v2 (CC BY 4.0)}
\begin{document}
\maketitle

\noindent\emph{Source and licence.} Late-time tails for linear waves on radially symmetric stationary spacetimes of two space dimensions. Version arXiv:2605.03220v2 (CC BY 4.0), distributed by arXiv under the Creative Commons Attribution 4.0 International licence, http://creativecommons.org/licenses/by/4.0/, as recorded in the arXiv licence metadata for this exact version and shown on the abstract page https://arxiv.org/abs/2605.03220v2. Full licence notice: \texttt{licenses/arxiv-2605.03220-CC-BY-4.0.txt}.\par\smallskip

\noindent\textit{Editorial scope.} Counted unit: the article's lemma data-bound (source0.tex lines 2149--2162). The article's complete original proof of it is delivered with it (source0.tex lines 2171--2230, one \texttt{proof} environment of 60 lines, which the article places away from the statement). No mathematics is written, repaired or re-derived: the statement and the proof are the article's own lines, in the article's own notation, and no retained line is substituted or re-flowed.  No further source-local context is needed: every cross-reference of every retained line resolves inside the two delivered spans.  Nothing was omitted from any retained span, so the package carries no omission record.  No retained line contains a citation, so the wrapper carries no bibliography and no bibliography entry.  No retained line contains a figure environment, an image-inclusion command or drawing code, so no image is delivered and the package declares no asset dependency.  Editorial formatting only, in addition: an AMS \texttt{amsart} preamble that restates the article's own declarations and macros used by the retained lines, namely the environments \texttt{lemma} on separate counters, together with the standard packages those lines need.  The retained environments print in this document as Lemma 1 in order of appearance, whereas the article prints the same items with the numbering of its own full document.  The complete native source of the article as distributed in the arXiv e-print for this exact version is bundled unchanged as \texttt{sources/199778/source0.tex}.
\begin{lemma}[Abstract interpolation lemma]
Let \(\alpha{},\beta{},\delta{},D\ge 0\), and let \(f:[0,\infty)\to
\R_{\ge 0}\) be an integrable function satisfying the following assumptions:
\begin{align}
f(0)&\lesssim D,\label{data-bound}\\
f(\tau{}')&\lesssim f(\tau{}) + (1+\tau{})^{-\beta{}}D\quad \textnormal{for }\tau{}'\ge \tau{}\ge 0, \label{energy-bound}\\
\int_{\tau{}}^{\infty} f(\tau{}')\dd{}\tau{}'&\lesssim R^{-\alpha{}}(1+\tau{})^\delta{}D + R(f(\tau{}) + (1+\tau{})^{-\beta{}}D)\quad \textnormal{for }1\le R\le \tau{}. \label{dyadic-bound}
\end{align}
Then for \(\tau{}\ge 0\) and \(\eta{} > 0\), we have
\begin{equation}\label{iteration-conc}
f(\tau{})\lesssim_{\alpha{},\beta{},\eta{}} (1+\tau{})^{-\min (\beta{},\alpha{}+1) + \delta{} + \eta{}}D.
\end{equation}
\label{interpolation}
\end{lemma}

\begin{proof}
By \cref{energy-bound} and \cref{data-bound}, we have
\begin{equation}\label{energy-bound-cons}
f(\tau{})\lesssim f(0) + D\lesssim D\quad \textnormal{for }\tau{}\ge 0,
\end{equation}
and so it suffices to prove \cref{iteration-conc} for \(\tau{}\ge 1\).

\step{Step 1: The iteration argument.} Suppose we
have shown that
\begin{equation}\label{iteration-1}
f(\tau{})\lesssim \tau^{-\gamma{}+\delta{}}D\quad \textnormal{for }\tau{}\ge 1
\end{equation}
for some \(0\le \gamma{}\le \min (\beta{},\alpha{}+1)\). We will show that \cref{iteration-1} holds with
\(\min (1 + (1-\frac{1}{\alpha{}+1})\gamma{},\beta{})\) in place of \(\gamma{}\)
(with an implicit constant that is larger by a factor independent of
\(\gamma{}\)). Note that the case \(\gamma{} = 0\) of \cref{iteration-1} follows
from \cref{energy-bound-cons}.

Let \(\tau_i=2^i\) be a dyadic sequence. Combine \cref{iteration-1} with
\cref{dyadic-bound} to obtain
\begin{equation}\label{iteration-2}
\int_{\tau_i}^{\tau_{i + 1}} f(\tau{})\dd{}\tau{}\lesssim R^{-\alpha{}}\tau_i^\delta{}D + R\tau_i^{-\gamma{}}D\quad \textnormal{for }1\le R\le \tau{}_i.
\end{equation}
Since \(\gamma{}\le \alpha{} + 1\), we can take \(R = \tau_i^{\gamma{}/(\alpha{}+1)}\) in \cref{iteration-2} to obtain
\begin{equation}
\int_{\tau_i}^{\tau_{i + 1}} f(\tau{})\dd{}\tau{}\lesssim \tau{}_i^{-\gamma{}\alpha{}/(\alpha{}+1)+\delta{}}D.
\end{equation}
By the pigeonhole principle, there is \(\tau_i'\in [\tau_i,\tau_{i + 1}]\) such that
\begin{equation}\label{iteration-3}
f(\tau{}_i')\lesssim \tau{}_i^{-1-\gamma{}\alpha{}/(\alpha{}+1)+\delta{}}D.
\end{equation}
By \cref{energy-bound} and the dyadicity of the \(\tau_i\), and hence of the \(\tau_i'\),
we can extend the estimate \cref{iteration-3} to
\begin{equation}\label{iteration-4}
f(\tau{})\lesssim \tau{}^{-\min (1+\gamma{}\alpha{}/(\alpha{}+1),\beta{})+\delta{}}D\quad \textnormal{for }\tau{}\ge 1,
\end{equation}
as desired.

\step{Step 2: Completing the proof.} Define a sequence \(\gamma_n\) for \(n\ge 0\) by
\begin{equation}
\gamma{}_0=0,\qquad \gamma{}_{n+1} = \min \Bigl(1 + \Bigl(1-\frac{1}{\alpha{}+1}\Bigr)\gamma{}_n,\beta{}\Bigr).
\end{equation}
If \(\gamma_n\le \alpha{} + 1\), then so is \(\gamma_{n + 1}\) (since \(\gamma_{n + 1}\le 1 + (1-1/(\alpha{} +
1))\gamma_n\)). Since \(\gamma_0 = 0\), it follows that \(\gamma_n\le \min (\beta{},\alpha{} +
1)\) for all \(n\ge 0\). In particular, Step 1 shows that
\begin{equation}
f(\tau{})\lesssim_n \tau^{-\gamma{}_n+\delta{}}D\quad \textnormal{for }\tau{}\ge 1.
\end{equation}
If \(\beta{}\ge \alpha{} + 1\), then by induction we have
\begin{equation}
\gamma{}_{n+1} = 1 + \Bigl(1-\frac{1}{\alpha{}+1}\Bigr)\gamma{}_n\implies \gamma{}_{n+1} = \sum_{k=0}^n\Bigl(1-\frac{1}{\alpha{}+1}\Bigr)^k = \alpha{}+1-\alpha{}\Bigl(1-\frac{1}{\alpha{}+1}\Bigr)^n.
\end{equation}
If \(\beta{}<\alpha{} + 1\), then we have
\begin{equation}
\gamma{}_{n+1} = \min \Bigl(\alpha{}+1-\alpha{}\Bigl(1-\frac{1}{\alpha{}+1}\Bigr)^n,\beta{}\Bigr),
\end{equation}
and there is some \(N\ge 0\) such that \(\gamma_n=\beta{}\) for \(n\ge N\). In either case, we
establish \cref{iteration-conc} by taking \(n\) sufficiently large that
\(\alpha{}(1-1/(\alpha{} + 1))^n \le \eta{}\).
\end{proof}

\end{document}
